\section{Integral cohomology}
Let $F$ be a smooth fibre, marked so that
\[
 H_1(F;\Z)=\Lambda,
 \qquad
 H^1(F;\Z)=V,
 \qquad
 H^q(F;\Z)=\bigwedge^qV.
\]
Let $j:B^\circ\hookrightarrow B$ and let $L^q$ be the local system
$R^qf_*\Z|_{B^\circ}$.

\subsection{Specialization at the multiple fibres}
Put
\[
 \psi_1=2u+w+3\delta,
 \qquad
 \psi_2=u+w+2\delta.
\]
The central torus in the covering model of $N_j$ maps to $S_j$ by an
$m_j$-sheeted covering.  Using the smooth marking over $\Delta_j$, identify
that torus with a nearby fibre $F$ and denote the covering by
$\pi_j:F\to S_j$.

\begin{proposition}\label{prop:finite-specialization}
The groups $H^q(S_j;\Z)$ are torsion-free of ranks
\[
 (1,2,2,2,1),
\]
and the pullback from $S_j$ to a nearby torus is injective.  Its image in
the monodromy invariants has the following indices:
\[
\begin{array}{c|ccccc}
q&0&1&2&3&4\\ \hline
j=1&1&3&1&1&3\\
j=2&1&4&2&2&4.
\end{array}
\]
More precisely,
\[
 \pi_1^*H^1(S_1)=\langle3\gamma,\psi_1\rangle,
 \qquad
 \pi_2^*H^1(S_2)=\langle4\gamma,\psi_2\rangle.
\]
The degree-two cokernel at $p_2$ is generated by the class of
$q=uw+6\gamma\delta$, and the degree-three cokernel at $p_2$ by the class
of $\gamma uw$.
\end{proposition}

\begin{proof}
The fundamental group of $S_j$ has the presentation
\[
 \langle\Lambda,\wh g_j\mid
 \wh g_j\lambda\wh g_j^{-1}=A_j\lambda,
 \ \wh g_j^{m_j}=t_{v_j}\rangle.
\]
The matrices in \cref{prop:linear-algebra} give
\[
 (A_1-I)\Lambda
 =\{\lambda:\gamma(\lambda)=0,
   \ \psi_1(\lambda)=0\},
 \qquad
 (A_2-I)\Lambda
 =\{\lambda:\gamma(\lambda)=0,
   \ \psi_2(\lambda)=0\}.
\]
Indeed the generators of $(A_1-I)\Lambda$ are
$6\wh u-6\wh w-2\wh\delta$,
$-\wh u-\wh w+\wh\delta$ and
$\wh u-2\wh w$, which generate the displayed kernel; the analogous
statement follows from the three generators of $(A_2-I)\Lambda$.
The maps $(\gamma,\psi_j):\Lambda\to\Z^2$ are surjective, so
$\Lambda/(A_j-I)\Lambda\cong\Z^2$.

After abelianization the relation is
\[
 m_j\wh g_j=(\gamma(v_j),\psi_j(v_j)).
\]
For $v_1=\eps$ this is $(1,0)$ and for $v_2=-\eps'$ it is $(-1,0)$.
The relation vector in $\Z^3$ is therefore primitive, and
$H_1(S_j;\Z)\cong\Z^2$ in both cases.  Since $S_j$ is a free finite
quotient of a four-torus, its Euler characteristic is zero.  Thus
$b_1=b_3=2$ and $b_2=2$.  Poincaré duality gives
$H^3(S_j;\Z)\cong H_1(S_j;\Z)\cong\Z^2$.  In the universal coefficient
sequence for $H^3$, the torsion group
$\operatorname{Ext}(H_2(S_j;\Z),\Z)$ injects into this free group; hence
$H_2(S_j;\Z)$ is torsion-free.  Poincaré duality and the universal
coefficient theorem now show that every $H^q(S_j;\Z)$ is torsion-free,
with ranks $(1,2,2,2,1)$.  The cohomological transfer of $\pi_j$ satisfies
$\operatorname{tr}\circ\pi_j^*=m_j$; hence $\pi_j^*$ is injective in every
degree.

A homomorphism $\pi_1(S_j)\to\Z$ restricting to
$\xi\in V^{T_j}$ exists exactly when $m_j$ divides $\xi(v_j)$.  Since
\[
 V^{T_1}=\langle\gamma,\psi_1\rangle,
 \qquad
 V^{T_2}=\langle\gamma,\psi_2\rangle,
\]
and $\psi_j(v_j)=0$, the degree-one images are the displayed lattices.

The invariant degree-two lattices are
\[
 (\textstyle\bigwedge^2V)^{T_1}=\langle\gamma\psi_1,q\rangle,
 \qquad
 (\textstyle\bigwedge^2V)^{T_2}=\langle\gamma\psi_2,q\rangle.
\]
Their intersection matrices, with values in $\Z\vol$, are
\[
 \begin{pmatrix}0&3\\3&12\end{pmatrix},
 \qquad
 \begin{pmatrix}0&2\\2&12\end{pmatrix},
\]
with determinants $-9$ and $-4$.  The intersection form on
$H^2(S_j;\Z)$ is unimodular.  It is even.  Indeed, the covering torus has
a nowhere-vanishing holomorphic two-form $\omega_0$.  At the central
point the linear part of the deck generator is $R_j(z_j)$, so
$\omega_0$ is multiplied by $\det R_j(z_j)$.  From
\cref{prop:period-equivariance,eq:mu-values},
\[
 \det R_1(z_1)=-\rho^{-1}=e^{2\pi i/3},
 \qquad
 \det R_2(z_2)=i^{-1}=-i.
\]
The form $\omega_0$ identifies the pullback of $K_{S_j}$ with the
trivial line bundle equipped with the character
$\wh g_j\mapsto\det R_j(z_j)^{-1}$.  This character has finite order, so
$K_{S_j}$ is torsion and hence numerically trivial.  Wu's formula gives
$x^2\equiv x\cdot c_1(K_{S_j})\equiv0\pmod2$ for every
$x\in H^2(S_j;\Z)$.  Pullback multiplies the intersection form by the
covering degree $m_j$.
If $k_j$ is the index of its image in the invariant lattice, determinant
comparison gives
\[
 k_1^2=\frac{3^2}{9}=1,
 \qquad
 k_2^2=\frac{4^2}{4}=4.
\]
Thus the degree-two indices are $1$ and $2$.  The vector $q$ is not in the
image for $j=2$: otherwise $q=\pi_2^*a$ would give
$12=q^2=4a^2$, so $a^2=3$, contradicting evenness.  Hence its class
generates the cokernel.

The degree-three indices follow from the perfect pairing
$H^1\times H^3\to\Z$.  Direct calculation gives
\[
 (\textstyle\bigwedge^3V)^{T_1}
 =\langle\gamma uw,-uw\delta+2\gamma w\delta+4\gamma u\delta\rangle,
\]
\[
 (\textstyle\bigwedge^3V)^{T_2}
 =\langle\gamma uw,-uw\delta+3\gamma w\delta+3\gamma u\delta\rangle.
\]
The determinants of the pairings with the degree-one invariant lattices
are $3$ and $2$.  Let $k_j^{(r)}$ be the index of the pullback image in
degree $r$.  If bases of the pullback lattices are expressed in bases of
the invariant lattices, the determinants of the two change-of-basis
matrices are $k_j^{(1)}$ and $k_j^{(3)}$.  On the other hand, pullback
multiplies the perfect pairing on $S_j$ by $m_j$, so the determinant of
the restricted rank-two pairing is $m_j^2$.  Therefore
\[
 k_j^{(1)}k_j^{(3)}
 \,|\det( (\textstyle\bigwedge^1V)^{T_j}
            \times(\textstyle\bigwedge^3V)^{T_j})|
 =m_j^2.
\]
Using $k_1^{(1)}=3$, $k_2^{(1)}=4$ gives
$k_1^{(3)}=1$, $k_2^{(3)}=2$.  Moreover $\psi_2$ lies in the degree-one
pullback image, whereas every pairing between pullback classes in degrees
one and three is divisible by $4$.  Since
$\langle\psi_2,\gamma uw\rangle=-2$, the class $\gamma uw$ is not in
the degree-three pullback image and therefore generates its order-two
cokernel.  In degree four pullback is multiplication by the covering
degree.
\end{proof}

\subsection{Specialization at the cusp}
\begin{theorem}[Integral invariant cycles at the cusp]\label{thm:cusp-specialization}
For every $q$, the specialization map
\[
 \operatorname{sp}^q:H^q(W;\Z)\longrightarrow H^q(F;\Z)
\]
is injective and has image
\[
 H^q(F;\Z)^{T_0}.
\]
In particular
\[
 H^q(W;\Z)\cong
 \Z,\ \Z^2,\ \Z^4,\ \Z^2,\ \Z
 \qquad(q=0,1,2,3,4),
\]
and all these groups are torsion-free.
\end{theorem}

\subsection{The integral direct-image sheaves}
Write
\[
 R^q=R^qf_*\Z.
\]
Proper base change identifies the stalk of $R^q$ at a special point with
the cohomology of the corresponding fibre.  At $p_0$, the map to the
nearby invariant lattice is the specialization map of
\cref{thm:cusp-specialization}.  At $p_j$, the local tube retracts to
$S_j$, and the map to a nearby fibre is the pullback along the covering
$F\to S_j$.  These maps are injective by
\cref{prop:finite-specialization,thm:cusp-specialization}; hence $R^q$ is
the subsheaf of $j_*L^q$ determined by the three displayed stalk lattices.

\begin{lemma}[Cohomology of the extended local systems]
\label{lem:local-system-cohomology}
Let $L$ be a local system of free finitely generated abelian groups on
$B^\circ$, with stalk $M$ and monodromy group $G$.
\begin{enumerate}[label=\textup{(\roman*)}]
\item There is a natural isomorphism
\begin{equation}\label{eq:H2-coinvariants}
 H^2(B,j_*L)\cong M_G.
\end{equation}
It is compatible with pairings of local systems: under
\eqref{eq:H2-coinvariants}, cup product by an invariant element is the
induced action on coinvariants.
\item Let $t_1,t_2$ be the clockwise generators of
$\pi_1(B^\circ)$, let $t_0=(t_1t_2)^{-1}$, and write $T_i$ for their
actions on $M$.  Then
\begin{equation}\label{eq:parabolic-cohomology}
 H^1(B,j_*L)
 \cong
 \frac{(T_1-I)M\cap(T_0-I)M}{(T_1-I)M^{T_2}}.
\end{equation}
\end{enumerate}
\end{lemma}

\begin{proof}
The quotient $j_*L/j_!L$ is supported at the three punctures and has no
cohomology in degree two.  Therefore
\[
 H^2(B,j_*L)\cong H_c^2(B^\circ,L)
 \cong H_0(B^\circ,L)\cong M_G
\]
by Poincaré duality with local coefficients \cite[Ch.~V]{Bredon}.
Naturality of cap product gives compatibility with coefficient pairings.

A class in $H^1(B,j_*L)$ may be represented by an inhomogeneous group
cocycle on $\pi_1(B^\circ)$.  Extension across the puncture $p_i$ means
that its restriction to the local cyclic subgroup is a coboundary; this
is equivalent to the condition $c(t_i)\in(T_i-I)M$.  Thus the relevant
classes are precisely the parabolic cocycles.  Subtracting a coboundary,
arrange $c(t_2)=0$; the remaining coboundaries are those defined by
elements of $M^{T_2}$.  Put
$x=c(t_1)$.  The cocycle identity and $t_0=(t_1t_2)^{-1}$ give
$c(t_0)=-T_0x$.  Hence the parabolic conditions are
$x\in(T_1-I)M\cap(T_0-I)M$, and the residual coboundaries change $x$ by
$(T_1-I)M^{T_2}$.  This proves \eqref{eq:parabolic-cohomology}.
\end{proof}

\begin{proposition}\label{prop:Rq-cohomology}
The quotient $j_*L^q/R^q$ is a skyscraper sheaf with stalks
\[
\begin{array}{c|ccc}
q&p_0&p_1&p_2\\ \hline
1&0&\Z/3&\Z/4\\
2&0&0&\Z/2\\
3&0&0&\Z/2\\
4&0&\Z/3&\Z/4.
\end{array}
\]
Consequently
\begin{align}
 H^0(B,R^1)&=12\Z\gamma,
 &H^0(B,R^2)&=2\Z q,
 \label{eq:H0R12}\\
 H^0(B,R^3)&=2\Z\gamma uw,
 &H^0(B,R^4)&=12\Z\vol.
 \label{eq:H0R34}
\end{align}
Furthermore
\begin{align}
 H^2(B,R^1)&=\Z[\delta],
 &H^2(B,R^2)&=\Z[\gamma\delta],
 \label{eq:H2R12}\\
 H^2(B,R^3)&=\Z[uw\delta],
 &H^2(B,R^4)&=\Z[\vol],
 \label{eq:H2R34}
\end{align}
and
\begin{equation}\label{eq:H1-vanishing}
 H^1(B,R^1)=H^1(B,R^2)=0.
\end{equation}
The identifications in degree two are compatible with cup products: an
invariant class acts on coinvariants by exterior product.
\end{proposition}

\begin{proof}
The skyscraper table is exactly
\cref{prop:finite-specialization,thm:cusp-specialization}.  Taking global
invariants and imposing the local divisibilities gives
\eqref{eq:H0R12}--\eqref{eq:H0R34}, using
\cref{prop:linear-algebra}(vi).

Since the skyscraper quotient $j_*L^q/R^q$ has no degree-two
cohomology, \cref{lem:local-system-cohomology}(i) and the coinvariant
calculations of \cref{prop:linear-algebra}(vi) give
\eqref{eq:H2R12}--\eqref{eq:H2R34} with the stated compatibility.

The maps from the global invariant lattices to the skyscraper quotients
are surjective in degrees one and two: in degree one, $\gamma$ maps to
$(1,1)\in\Z/3\oplus\Z/4$, and in degree two, $q$ maps to the generator of
$\Z/2$.  Hence $H^1(B,R^q)=H^1(B,j_*L^q)$ for $q=1,2$.  Apply
\cref{lem:local-system-cohomology}(ii).  For $M=V$,
\[
 (T_1-I)V=\langle2\gamma+u+w,2\gamma-u\rangle,
 \qquad
 (T_0-I)V=\langle\gamma,u\rangle,
\]
whose intersection is $\Z(2\gamma-u)$.  Since
$(T_1-I)\psi_2=-(2\gamma-u)$, this equals $(T_1-I)V^{T_2}$.  For
$M=\bigwedge^2V$,
\[
 (\textstyle\bigwedge^2T_0-I)M
 =\langle\gamma u,\gamma w+u\delta\rangle,
\]
while $(\bigwedge^2T_1-I)M$ is generated by
\[
 \gamma u,\ \gamma w,\ 2u\delta+w\delta,
 \ uw+u\delta-w\delta-6\gamma\delta.
\]
The intersection is $\Z\gamma u$.  Since $\gamma\psi_2$ is fixed by
$\bigwedge^2T_2$ and
$(\bigwedge^2T_1-I)(\gamma\psi_2)=\gamma u$, this is
$(\bigwedge^2T_1-I)M^{T_2}$.  Formula
\eqref{eq:parabolic-cohomology} therefore gives the required vanishing.
\end{proof}

\subsection{The Leray differential}
Consider the integral Leray spectral sequence
\[
 E_2^{p,q}=H^p(B,R^qf_*\Z)\Longrightarrow H^{p+q}(X;\Z).
\]
Since $B$ has real dimension two, the only possible nonzero differentials
are
\[
 d_2^{0,q}:E_2^{0,q}\longrightarrow E_2^{2,q-1}.
\]

\begin{proposition}\label{prop:leray-differentials}
After choosing the sign of the generator
$\omega\in H^2(B;\Z)$, one has
\begin{align}
 d_2(12\gamma)&=\omega,
 \label{eq:d2-gamma}\\
 d_2(2q)&=\omega[\delta],
 \label{eq:d2-q}\\
 d_2(2\gamma uw)&=\omega[\gamma\delta].
 \label{eq:d2-gamma-uw}
\end{align}
In particular all three maps are isomorphisms of infinite cyclic groups.
\end{proposition}

\begin{proof}
The first differential is the obstruction to lifting the global section
$12\gamma$ of $R^1$ to a class in $H^1(X;\Z)$.  We compute it by local
lifts.  Let $U=D_0\sqcup D_1\sqcup D_2$ be a disjoint union of small
discs and let $V$ be a neighbourhood of the complementary pair of pants,
so that $U\cap V$ is a disjoint union of three annuli.  For each of the three discs and for $V$, the five-term Leray exact
sequence contains
\[
 H^1(f^{-1}(Y);\Z)\longrightarrow H^0(Y,R^1)
 \longrightarrow H^2(Y;\Z),
\]
where $Y$ denotes the corresponding base region.  The last group vanishes.
Thus a section of $R^1$ lifts to a degree-one class on the inverse image
exactly when it lies in the appropriate local specialization lattice.
This holds for $12\gamma$ by
\cref{prop:finite-specialization,thm:cusp-specialization}.
On each annulus the two lifts have the same image in $H^0(R^1)$.  The
five-term exact sequence for the annulus therefore identifies their
difference with the pullback of a unique class in $H^1$ of the annulus.
These three classes form a Čech cocycle, and its image under the
Mayer--Vietoris isomorphism
\[
 H^1(U\cap V;\Z)/\im H^1(V;\Z)
 \xrightarrow{\sim}H^2(B;\Z)
\]
is exactly $d_2(12\gamma)$.  Evaluating a difference on the oriented
annular meridian identifies $H^1(U\cap V;\Z)$ with $\Z^3$; the image of
$H^1(V;\Z)$ is the sum-zero sublattice, so the class of
$(n_0,n_1,n_2)$ is $\pm(n_0+n_1+n_2)\omega$.

Over $D_0$, choose the lift to have value zero on the toric meridian,
which is null-homotopic in $N_0$.  Over $D_j$, if $\Theta_j$ is a lift,
the relation $\wh g_j^{m_j}=t_{v_j}$ gives
\[
 m_j\Theta_j(\wh g_j)=12\gamma(v_j).
\]
Thus the meridian values of the local lifts are
\[
 4\quad\text{at }p_1,
 \qquad
 -3\quad\text{at }p_2,
 \qquad
 0\quad\text{at }p_0.
\]
Let the lift over $V$ have values $a_1,a_2$ on the first two meridians.
The product relation for the three clockwise boundary meridians gives
value $-a_1-a_2$ at the cusp.  The three differences between the disc
lifts and the $V$-lift are therefore
\[
 4-a_1,
 \qquad
 -3-a_2,
 \qquad
 a_1+a_2.
\]
Their sum is $1$.  By the preceding Mayer--Vietoris identification,
choosing the sign of $\omega$ accordingly proves \eqref{eq:d2-gamma}.

Write
\[
 d_2(2q)=a\,\omega[\delta],
 \qquad
 d_2(2\gamma uw)=b\,\omega[\gamma\delta]
\]
for integers $a,b$.  Multiplicativity and the relation
\[
 (12\gamma)(2q)=12(2\gamma uw)
\]
give, using the graded Leibniz rule with the sign $(-1)^1=-1$ and the
identity $[q]=12[\gamma\delta]$ in the degree-two coinvariants,
\[
 12b=24-12a,
 \qquad\text{hence}\qquad b=2-a.
\]
Next $(2q)(2\gamma uw)=0$ because $H^5(F;\Z)=0$.  Under the
coinvariant identifications of \cref{lem:local-system-cohomology}, the two
coefficient pairings
\[
 H^2(B,R^1)\otimes H^0(B,R^3)\longrightarrow H^2(B,R^4),
 \qquad
 H^0(B,R^2)\otimes H^2(B,R^2)\longrightarrow H^2(B,R^4)
\]
have absolute value $2$ on the displayed generators: on representatives,
\[
 \delta\wedge(2\gamma uw)=-2\vol,
 \qquad
 (2q)\wedge(\gamma\delta)=2\vol.
\]
The graded Leibniz rule applied to the zero product therefore gives
$|a|=|b|$.  Since $b=2-a$, one has $a^2=(2-a)^2$, hence $a=b=1$.
This proves
\eqref{eq:d2-q}--\eqref{eq:d2-gamma-uw}.
\end{proof}

\begin{theorem}\label{thm:integral-cohomology}
The integral homology of $X$ is that of $S^6$.
\end{theorem}

\begin{proof}
By \cref{prop:Rq-cohomology,prop:leray-differentials}, the spectral
sequence gives
\[
 H^1(X;\Z)=H^2(X;\Z)=H^3(X;\Z)=0.
\]
Indeed, in total degree one the only term is the kernel of
\eqref{eq:d2-gamma}; in degree two the three graded pieces are the
cokernel of \eqref{eq:d2-gamma}, $H^1(B,R^1)$, and the kernel of
\eqref{eq:d2-q}; in degree three they are the cokernel of
\eqref{eq:d2-q}, $H^1(B,R^2)$, and the kernel of
\eqref{eq:d2-gamma-uw}.

Since $\pi_1(X)=1$ by \cref{thm:pi1}, one has $H_1(X;\Z)=0$.
The universal coefficient theorem and $H^2(X;\Z)=0$ imply
$\Hom(H_2(X),\Z)=0$, so the finitely generated group $H_2(X)$ is torsion.
The group $\operatorname{Ext}(H_2(X),\Z)$ injects into $H^3(X;\Z)=0$;
for a finite abelian group it is isomorphic to the group itself.  Hence
$H_2(X;\Z)=0$.  Integral Poincaré duality now gives
\[
 H_3(X;\Z)=H_4(X;\Z)=H_5(X;\Z)=0.
\]
Finally $H_0(X;\Z)=H_6(X;\Z)=\Z$ because $X$ is connected and oriented.
\end{proof}

